\section{Introduction}
\label{sec:introduction}
\IEEEPARstart{G}{enerative} design is increasingly used to explore large engineering design spaces, while aerospace workflows continue to rely heavily on computational fluid dynamics (CFD), surrogate models, and optimization loops to evaluate aerodynamic tradeoffs \cite{elham2021openfemflow, andres2021surrogate_aero}. Aircraft design adds further constraints beyond raw shape generation, including manufacturability, structural load paths, and mission-specific performance targets. Those requirements make it easy to overstate what a proof-of-concept generative repository can presently demonstrate.

This paper therefore takes a deliberately narrower position. The repository implements a latent diffusion-style pipeline that maps latent samples to voxel geometries, scores them with an internal CFD-oriented path, and exports them for downstream validation. The present evidence comes from synthetic training data and reduced sanity experiments, so the goal of this paper is not to claim a finished aircraft-design system. Instead, the goal is to document the implementation, report what the current code path does reproduce, and make the boundary between demonstrated behavior and future work explicit.

The key contributions of this work are:
\begin{enumerate}
    \item A reproducible proof-of-concept latent generative pipeline for producing aircraft-like voxel geometries and exporting them to mesh-based artifacts.
    \item An implementation path for CFD-informed evaluation that couples an internal lattice-Boltzmann-based scorer with an external OpenFOAM benchmark route, together with a reduced sanity experiment showing that the end-to-end path executes.
    \item An issue-driven validation framework that distinguishes code-path validation from stronger claims such as aircraft-specific validity, aerodynamic optimization, structural viability, and conditioned generation from mission or manufacturing constraints.
\end{enumerate}

The rest of the paper is organized as follows. Section \ref{sec:related_work} positions the repository relative to 3D generative modeling, topology optimization, and CFD-driven design practice. Section \ref{sec:methodology} describes the implementation that currently exists in the codebase. Section \ref{sec:results} reports the reduced sanity-run evidence and explains its limits. Section \ref{sec:validation} records the validation requirements that remain before stronger claims can be made. Section \ref{sec:conclusion} summarizes the present scope and the minimum future work needed for a conditioned airplane generator.
